\begin{figure}[htbp]
\centering
\includegraphics[width=0.95\linewidth]{crt_syntax.pdf}
\caption{TOC 現状ツリー（CRT）における因果結合の基本パターン}
\label{fig:crt_syntax}
\end{figure}

\begin{enumerate}
    \item \textbf{単一因果（Direct Cause）}：
    「もし原因Aが存在するならば、必然的に結果Bが生じる」という直接的な十分条件の結合である。単一の入力ノイズや物理的トリガーが、直接的な生体反射や状態遷移を確定させる関係を記述する。
    
    \item \textbf{AND結合（積結合 / Sufficient Cause）}：
    原因Aと原因Bが\textbf{同時に揃ったときにのみ}、結果Cが生じる結合である。TOC では伝統的に因果線を束ねる楕円（ANDノード）として表現される。
    
    臨床において極めて重要なのがこの結合である。例えば、「病床の父を看病しなければならないという規範」単独ではヒステリー麻痺は起きず、「自立して自由になりたいという生存欲求」単独でも麻痺は起きない。両者が同時に同一のプロセッサ内で実行され、互いに譲らないデッドロック（相矛盾する2つの要請）を形成した瞬間にのみ、システムは異常停止（身体化への迂回）を引き起こす。
    
    \item \textbf{OR結合（独立和 / Independent Causes）}：
    原因Aまたは原因Bのいずれか一方が単独で存在するだけで、結果Cが独立して引き起こされる結合である。複数の異なる外的刺激が、同一の情動的インターラプション（$\Lambda_{\mathrm{fear}}$ 等）を惹起する経路を記述する。
    
    \item \textbf{エンティティの階層トポロジー}：
    CRT は、下層から上層に向かって以下の3層構造を形成する。
    \begin{itemize}
        \item \textbf{根本原因（Root Cause / 緑色ノード）}：患者の環境、不可避な外的制約、あるいは生物学的本能などの基底前提。
        \item \textbf{中間エンティティ（黄色ノード）}：根本原因の衝突によって引き起こされた、内部ステートの未解決ループ（$\Omega_{\mathrm{cause}} = \emptyset |$）や情動スタックの滞留。
        \item \textbf{望ましくない結果（UDE / 赤色ノード）}：最上層に出力された、客観的に観測可能なハードウェアフォールト（痙性麻痺、失声、激しい咳などの身体化症状）。
    \end{itemize}
\end{enumerate}

